\documentclass{article}
\usepackage[utf8]{inputenc}
\usepackage[russian]{babel}
\usepackage{amsmath}
\usepackage{amssymb}
\usepackage{geometry}
\usepackage{pgfplots}
\geometry{a4paper, margin=1in}

\title{Лабораторная работа №4}
\author{}
\date{}

\begin{document}

\maketitle

\section{Контекстная свобода}
Докажем, что порождаемый язык не является контекстно-свободным. Рассмотрим его пересечение с регулярным языком \(a^*\). Тогда в выводе могут использоваться только правила 1 и 3, так как другие правила содержат символ \(b\) в правой части. Получаем язык слов вида \(a^{2^n - 2}\) для \(n > 1\). Известно, что язык \(\{a^{2^n - 2} \mid n > 1\}\) не является контекстно-свободным. Поскольку класс КС-языков замкнут относительно пересечения с регулярными языками, исходный язык также не может быть контекстно-свободным.

\section{Оптимизированный парсер}

В обычном наивном парсере при переборе возможных разбиений строки на поддеревья \(S\) выполняется слишком много проверок, что приводит к высокой вычислительной сложности. Для оптимизации заметим два важных инварианта, выполняющихся для любого слова, порождаемого нетерминалом \(S\):
\begin{enumerate}
    \item Слово всегда содержит подстроку ``aa''.
    \item Количество символов \(a\) в слове четно.
\end{enumerate}

При переборе возможных границ разбиения будем предварительно проверять эти два инварианта для выделяемой подстроки \(S\). Если инварианты не выполняются, немедленно переходим к следующему варианту разбиения, что позволяет значительно сократить пространство перебора.

Дополнительно введем мемоизацию вычислений, чтобы не работать с одной и той же подстрокой повторно.

В результате описанных оптимизаций поиск допустимого разбиения (если оно существует) происходит существенно быстрее, а вычислительная сложность алгоритма снижается.

\section{Оценка вычислительной сложности}

В оптимизированном парсере наиболее затратным является правило $S \rightarrow S S b S S$:

\begin{itemize}
    \item Внешний цикл перебирает $O(n)$ позиций символа 'b'
    \item Для каждой позиции выполняются два вложенных цикла:
    \begin{itemize}
        \item Перебор $O(n)$ способов разделения левой части на $S_1$ и $S_2$
        \item Перебор $O(n)$ способов разделения правой части на $S_3$ и $S_4$
    \end{itemize}
    \item Четыре рекурсивных вызова для подстрок: $S(k_1)$, $S(k_2)$, $S(k_3)$, $S(k_4)$
\end{itemize}

Верхняя граница сложности определяется произведением:
\[
O(n) \times O(n) \times O(n) \times O(n) = O(n^4)
\]
где:
\begin{itemize}
    \item Первый множитель $O(n)$ — перебор позиций 'b' во внешнем цикле
    \item Второй множитель $O(n)$ — перебор разделений левой части
    \item Третий множитель $O(n)$ — перебор разделений правой части
    \item Четвёртый множитель $O(n)$ — суммарный размер рекурсивных вызовов
\end{itemize}

С учётом мемоизации (каждая уникальная подстрока вычисляется один раз) итоговая сложность составляет $O(n^4)$ в худшем случае.

\section{Тестирование}

Для проверки эффективности предложенного алгоритма было проведено тестирование на двух типах данных:
\begin{itemize}
    \item словах, принадлежащих языку;
    \item словах, не принадлежащих языку.
\end{itemize}

Для каждой длины измерялось время выполнения двух алгоритмов:
\begin{itemize}
    \item обычного парсера;
    \item оптимизированного парсера, описанного в разделе 2.
\end{itemize}

Результаты представлены на графиках ниже. (Значения времени приведены в миллисекундах.)

\begin{figure}[h!]
\centering
\begin{minipage}{0.48\textwidth}
\centering
\begin{tikzpicture}
\begin{axis}[
    title={Обычный парсер: слова из языка},
    xlabel={Длина строки $n$},
    ylabel={Время $t$, мс},
    xmin=0, xmax=30,
    ymin=0, ymax=3300,
    xtick={10,15, 20, 25, 30},
    ytick={0.1,1, 20, 50, 3200},
    legend style={at={(0.5,-0.2)}, anchor=north},
    grid=both,
    width=\textwidth,
    height=0.8\textwidth
]

% Обычный парсер для слов из языка (синий)
\addplot[
    color=blue,
    mark=triangle,
    mark size=3pt,
    thick
]
coordinates {
    (9, 0.14)(15, 1.1)(20, 26.06)(24, 35.08)(30, 3251.33)
};
\addlegendentry{Обычный парсер}

\end{axis}
\end{tikzpicture}
\caption{Время разбора для слов из языка (обычный парсер)}
\label{fig:ordinary_in_lang}
\end{minipage}
\hfill
\begin{minipage}{0.48\textwidth}
\centering
\begin{tikzpicture}
\begin{axis}[
    title={Оптимизированный парсер: слова из языка},
    xlabel={Длина строки $n$},
    ylabel={Время $t$, мс},
    xmin=0, xmax=30,
    ymin=0, ymax=3,
    xtick={10,15,20, 25, 30},
    ytick={0,0.5,1,2,2.5},
    legend style={at={(0.5,-0.2)}, anchor=north},
    grid=both,
    width=\textwidth,
    height=0.8\textwidth
]

% Оптимизированный парсер для слов из языка (красный)
\addplot[
    color=red,
    mark=square,
    mark size=3pt,
    thick
]
coordinates {
    (9, 0.59)(15, 0.21)(20, 1.8)(24, 2.01)(30, 2.28)
};
\addlegendentry{Оптимизированный парсер}

\end{axis}
\end{tikzpicture}
\caption{Время разбора для слов из языка (оптимизированный парсер)}
\label{fig:optimized_in_lang}
\end{minipage}
\end{figure}

\begin{figure}[h!]
\centering
\begin{minipage}{0.48\textwidth}
\centering
\begin{tikzpicture}
\begin{axis}[
    title={Обычный парсер: слова не из языка},
    xlabel={Длина строки $n$},
    ylabel={Время $t$, мс},
    xmin=0, xmax=30,
    ymin=0, ymax=3300,
    xtick={10,15,20,25,30},
    ytick={1,10,100,300,3000},
    legend style={at={(0.5,-0.2)}, anchor=north},
    grid=both,
    width=\textwidth,
    height=0.8\textwidth
]

% Обычный парсер для слов не из языка (синий)
\addplot[
    color=blue,
    mark=triangle,
    mark size=3pt,
    thick
]
coordinates {
    (10, 0.12)(15, 0.86)(20, 14.19)(25, 335.54)(30, 3251.33)
};
\addlegendentry{Обычный парсер}

\end{axis}
\end{tikzpicture}
\caption{Время разбора для слов не из языка (обычный парсер)}
\label{fig:ordinary_not_in_lang}
\end{minipage}
\hfill
\begin{minipage}{0.48\textwidth}
\centering
\begin{tikzpicture}
\begin{axis}[
    title={Оптимизированный парсер: слова не из языка},
    xlabel={Длина строки $n$},
    ylabel={Время $t$, мс},
    xmin=0, xmax=30,
    ymin=0, ymax=3,
    xtick={10,15,20,25,30},
    ytick={0.05,0.1,0.5,1,2},
    legend style={at={(0.5,-0.2)}, anchor=north},
    legend style={at={(0.5,-0.2)}, anchor=north},
    grid=both,
    width=\textwidth,
    height=0.8\textwidth
]

% Оптимизированный парсер для слов не из языка (красный)
\addplot[
    color=red,
    mark=square,
    mark size=3pt,
    thick
]
coordinates {
    (10, 0.02)(15, 0.05)(20, 0.08)(25, 0.17)(30, 2.28)
};
\addlegendentry{Оптимизированный парсер}

\end{axis}
\end{tikzpicture}
\caption{Время разбора для слов не из языка (оптимизированный парсер)}
\label{fig:optimized_not_in_lang}
\end{minipage}
\end{figure}

\end{document}